The design of algorithms has traditionally focused on optimizing
resources like time, space, or communication 
without necessarily considering how hard it is to independently
certify correctness of results.  Yet, the need for certified or
auditable results is becoming increasingly important in a world where
untrusted components are unavoidable in the implementation of any
non-trivial software.

In this paper, we presented a principled
approach towards addressing this problem by showing that appropriate
measures of audit complexity can be treated as a first-class
computational resource to be optimized at the design stage of an
algorithm itself, rather than as an afterthought.  For deterministic
approximate counting, this approach yields completely new algorithms
that are sub-optimal with respect to time/space resources, and yet
permit significantly simpler certification in practice.  We believe
the same ``trade off'' applies to other hard combinatorial problems as
well, and one needs to navigate the algorithm design space carefully
to strike the right balance between audit complexity and
time/space/communication complexity.  We hope our work opens the door
for further such investigations.  As concrete steps for future work,
we wish to explore the audit complexity versus time/space complexity
tradeoff in probabilistic approximate counting.

\paragraph*{Acknowledgments}
This work was supported in part by National Research Foundation Singapore under its NRF Fellowship Programme[NRF-NRFFAI1-2019-0004 ] , Ministry of Education Singapore Tier 2 grant MOE-T2EP20121-0011, and Ministry of Education Singapore Tier 1 Grant [R-252-000-B59-114 ].

